\section*{Bab \ref{ch:g12:arith} --- Aritmetika}

\begin{solution}{exo:g12:arith:1}
$17 \times 119 = 2023$, jadi $2026 = 17 \times 119 + 3$: hasil baginya $119$,
sisanya $3$. Untuk $-2026$: $-2026 = 17\times(-120) + 14$ (memang
$17 \times 120 = 2040$ dan $2040 - 2026 = 14$): hasil baginya $-120$, sisanya
$14$ (sisanya harus terletak di $\intco{0}{17}$, jadi \emph{bukan} $-3$).
\end{solution}

\begin{solution}{exo:g12:arith:2}
Modulo $10$: $7^2 = 49 \equiv 9 \equiv -1$. Karena itu
$7^{100} = \left(7^2\right)^{50} \equiv (-1)^{50} = 1 \pmod{10}$: jadi angka
terakhir $7^{100}$ adalah $1$.
\end{solution}

\begin{solution}{exo:g12:arith:3}
Euklides: $1071 = 2\times462 + 147$; $462 = 3\times147 + 21$;
$147 = 7\times21 + 0$. Jadi $\gcd = 21$.

Substitusi mundurnya: $21 = 462 - 3\times147 = 462 - 3(1071 - 2\times462)
= 7\times462 - 3\times1071$. Jadi $u = -3$, $v = 7$:
$1071\times(-3) + 462\times7 = 21$.
\end{solution}

\begin{solution}{exo:g12:arith:4}
Setiap \omterm{def:g10:numbers:sets}{bilangan bulat} $\equiv 0, 1, 2$ atau $3 \pmod 4$, dan setelah dikuadratkan:
$0^2 \equiv 0$, $1^2 \equiv 1$, $2^2 = 4 \equiv 0$, $3^2 = 9 \equiv 1$. Jadi
$n^2 \equiv 0$ atau $1 \pmod 4$. Jumlah dua kuadrat lalu kongruen dengan
$0 + 0$, $0 + 1$ atau $1 + 1$, \emph{yaitu} dengan $0$, $1$ atau $2 \pmod 4$ ---
tak pernah dengan $3$.
\end{solution}

\begin{solution}{exo:g12:arith:5}
\omterm{def:g12:arith:divides}{Keterbagian} oleh $2$: di antara $n$ dan $n + 1$, salah satunya genap. \omterm{def:g12:arith:divides}{Keterbagian}
oleh $3$: jika $n \equiv 0$, maka $3 \mid n$; jika $n \equiv 1 \pmod 3$, maka
$2n + 1 \equiv 3 \equiv 0$; jika $n \equiv 2$, maka $n + 1 \equiv 0$. Pada semua
kasusnya $3$ \omterm{def:g12:arith:divides}{membagi} hasil kalinya. Karena $\gcd(2,3) = 1$,
\cref{cor:g12:arith:coprimeprod} memberi $6 \mid n(n+1)(2n+1)$.
(Ini sekaligus membuktikan kembali bahwa $\frac{n(n+1)(2n+1)}{6}$, jumlah kuadrat pada
\cref{exo:g12:seq:1}, adalah \omterm{def:g10:numbers:sets}{bilangan bulat}.)
\end{solution}

\begin{solution}{exo:g12:arith:6}
Kita mencari balikan $5$ modulo $11$: lewat pengujian (atau Bézout),
$5 \times 9 = 45 = 44 + 1 \equiv 1 \pmod{11}$. Dengan mengalikan \omterm{def:g12:arith:congruence}{kekongruenannya}
dengan $9$:
\[
x \equiv 9 \times 3 = 27 \equiv 5 \pmod{11}.
\]
Penyelesaiannya adalah \omterm{def:g10:numbers:sets}{bilangan bulat} $x = 5 + 11k$, $k \in \Z$. (Periksalah:
$5\times5 = 25 \equiv 3 \pmod{11}$.)
\end{solution}

\begin{solution}{exo:g12:arith:7}
$\gcd(17, 40) = 1$, jadi penyelesaiannya ada. Euklides:
$40 = 2\times17 + 6$; $17 = 2\times6 + 5$; $6 = 5 + 1$. Dengan substitusi mundur:
$1 = 6 - 5 = 6 - (17 - 2\times6) = 3\times6 - 17
= 3(40 - 2\times17) - 17 = 3\times40 - 7\times17$. Jadi
$17\times(-7) - 40\times(-3) = 1$: itulah penyelesaian khususnya,
$(x_0, y_0) = (-7, -3)$.

Penyelesaian umum $17x - 40y = 1$: dengan mengurangkan hubungan khususnya,
$17(x + 7) = 40(y + 3)$; karena $\gcd(17, 40) = 1$, Gauss memberi
$40 \mid x + 7$, sehingga $x = -7 + 40k$ lalu $y = -3 + 17k$, $k \in \Z$
(dan semuanya lolos pemeriksaan).

Untuk $17x - 40y = 6$, kalikan penyelesaian khususnya dengan $6$:
$(x_1, y_1) = (-42, -18)$, lalu penalaran yang sama memberi
\[
x = -42 + 40k, \qquad y = -18 + 17k, \qquad k \in \Z .
\]
(Misalnya $k = 2$: $x = 38$, $y = 16$; memang $17\times38 - 40\times16
= 646 - 640 = 6$.)
\end{solution}

\begin{solution}{exo:g12:arith:8}
Andaikan $\sqrt2 = \frac ab$ dengan $a, b \in \N^*$; maka $a^2 = 2b^2$. Pada
pemfaktoran \omterm{def:g12:arith:prime}{prima} suatu kuadrat, setiap eksponennya genap; jadi eksponen
$2$ pada $a^2$ genap, sedangkan pada $2b^2$ ganjil (satu lebih daripada bilangan
genap). Dua pemfaktoran \omterm{def:g10:numbers:sets}{bilangan bulat} yang sama dengan eksponen $2$ yang berbeda
bertentangan dengan ketunggalan pada \cref{thm:g12:arith:fta}. Jadi pecahan
semacam itu tak ada: $\sqrt2 \notin \Q$.
\end{solution}

\begin{solution}{exo:g12:arith:9}
\emph{1.} Dari $k\binom pk = p \binom{p-1}{k-1}$, jelas $p$ \omterm{def:g12:arith:divides}{membagi}
$k\binom pk$. Untuk $1 \leq k \leq p-1$, $p \nmid k$ dan $p$ yang \omterm{def:g12:arith:prime}{prima} memberi
$\gcd(p, k) = 1$, sehingga lema Gauss menghasilkan $p \mid \binom pk$.

\emph{2.} Induksi pada $a$. Untuk $a = 0$: $0^p \equiv 0$. Andaikan
$a^p \equiv a \pmod p$. Menurut teorema binomial,
\[
(a+1)^p = \sum_{k=0}^{p} \binom pk a^k
\equiv a^p + 1 \pmod p,
\]
sebab semua suku tengahnya lenyap modulo $p$ menurut butir 1. Menurut hipotesis
induksinya, $(a+1)^p \equiv a + 1 \pmod p$. Ini membuktikan
$a^p \equiv a$ untuk semua $a \in \N$, dan kasus $a < 0$ menyusul dengan menulis
$a \equiv a + kp$ untuk suatu wakil positif yang sesuai.
\end{solution}

\begin{solution}{exo:g12:arith:10}
$n \equiv 2 \pmod 3$ dan $n \equiv 3 \pmod 5$: tulislah $n = 2 + 3s$; maka
$2 + 3s \equiv 3 \pmod 5$, \emph{yaitu} $3s \equiv 1 \pmod 5$. Balikan
$3$ modulo $5$ adalah $2$ (sebab $3\times2 = 6 \equiv 1$), jadi $s \equiv 2 \pmod 5$,
sebutlah $s = 2 + 5t$, dan $n = 8 + 15t$: jadi kedua syarat pertamanya berarti
$n \equiv 8 \pmod{15}$.

Dengan menambahkan $n \equiv 2 \pmod 7$: $8 + 15t \equiv 2 \pmod 7$, dan
$15 \equiv 1 \pmod 7$, jadi $t \equiv -6 \equiv 1 \pmod 7$, sebutlah
$t = 1 + 7u$. Karena itu $n = 23 + 105u$:
\[
n \equiv 23 \pmod{105}.
\]
(Periksalah: $23 = 3\times7 + 2 = 5\times4 + 3 = 7\times3 + 2$.)
\end{solution}

\begin{solution}{pb:g12:arith:1}
\textbf{1.} $2026 = 289 \times 7 + 3$, jadi $2026 \equiv 3
\pmod 7$. Pangkat $7$ modulo $10$: $7, 9, 3, 1$, sebuah kitaran
sepanjang $4$; dan $100 \equiv 0 \pmod 4$, jadi angka terakhir
$7^{100}$ adalah $1$.

\textbf{2.} $5^2 = 25 \equiv -1 \pmod{13}$, sehingga
$5^{116} = \left(5^2\right)^{58} \equiv (-1)^{58} = 1$ dan
$5^{117} \equiv 5 \pmod{13}$.

\textbf{3.} Balikan $3$ modulo $7$ adalah $5$
(sebab $15 \equiv 1$): jadi $x \equiv 5 \times 5 = 25 \equiv 4 \pmod 7$.

\textbf{4.} $97 = 2 \times 35 + 27$; $35 = 27 + 8$;
$27 = 3 \times 8 + 3$; $8 = 2 \times 3 + 2$; $3 = 2 + 1$.
Dengan substitusi mundur: $1 = 97 \times 13 + 35 \times (-36)$. Jadi
$35 \times (-36) \equiv 1 \pmod{97}$: balikan $35$ adalah
$-36 \equiv 61 \pmod{97}$.

\textbf{5.} Bilangan $a$ punya balikan modulo $n$ tepat ketika
$\gcd(a, n) = 1$: Bézout menyediakan $au + nv = 1$, yaitu
$au \equiv 1$; sebaliknya adanya balikan memaksa FPB-nya \omterm{def:g12:arith:divides}{membagi}
$1$.

\textbf{6.} $0{\cdot}10 + 3{\cdot}9 + 0{\cdot}8 + 6{\cdot}7 +
4{\cdot}6 + 0{\cdot}5 + 6{\cdot}4 + 1{\cdot}3 + 5{\cdot}2 +
2{\cdot}1 = 132 = 12 \times 11 \equiv 0 \pmod{11}$: jadi sah.

\textbf{7.} Jumlahnya berubah sebesar $wd$ dengan
$1 \leq w \leq 10$ dan $1 \leq \abs d \leq 9$: karena $11$
\omterm{def:g12:arith:prime}{prima} dan tidak \omterm{def:g12:arith:divides}{membagi} satu pun faktornya, ia tak dapat \omterm{def:g12:arith:divides}{membagi} hasil kalinya
(\cref{thm:g12:arith:gauss} / \cref{prop:g12:arith:primedivides}),
sehingga jumlah yang berubah itu tak pernah lagi $\equiv 0$: jadi setiap
galat satu angka membunyikan alarmnya.

\textbf{8.} Menukar angka bertetangga $a, b$ (yang berbobot
$w + 1, w$) mengubah jumlahnya sebesar
$(w+1)b + wa - (w+1)a - wb = b - a \not\equiv 0$ bagi
$a \neq b$: jadi terdeteksi. Rahasianya adalah \emph{keprimaan}
$11$: modulo $10$, hasil kali seperti $5 \times 2$ lenyap padahal
tak satu faktornya nol, sehingga galat berbobot $5$ sebesar $\pm 2$ (atau
pertukaran yang bernasib sial) bisa lolos.

\textbf{9.} Jumlah berbobot kedua belas angkanya: $119$; angka
pemeriksanya harus melengkapinya menjadi kelipatan $10$, jadi $1$ (kode
lengkapnya $978\,2940199\,051$). EAN meleset pada pertukaran bertetangga
dengan $2(a - b) \equiv 0 \pmod{10}$, yaitu
$\abs{a - b} = 5$: menukar $2$ dengan $7$, misalnya, lolos
tanpa terlihat --- itulah harga \omterm{def:g12:complex:modulus}{modulus} $10$ yang bersahabat itu.

\textbf{10.} Dengan melipatduakan tiap angka kedua dari kanan lalu
melipatnya ($16 \to 7$, dan seterusnya), jumlahnya menjadi $80 \equiv 0
\pmod{10}$: jadi kartu ujinya sah.

\textbf{11.} Dengan \omterm{def:g12:complex:modulus}{modulus} \omterm{def:g12:arith:prime}{prima} setiap bobotnya punya balikan,
sehingga \emph{semua} galat tunggal dan \emph{semua} pertukaran
bertetangga tertangkap --- itulah kemewahan ISBN; sedangkan rancangan modulo
$10$ mempertahankan angka yang ramah bagi manusia dan menerima satu titik buta.

\textbf{12.} $2^{10} = 1024 = 3 \times 341 + 1 \equiv 1
\pmod{341}$, sehingga $2^{340} = \left(2^{10}\right)^{34} \equiv
1$. Padahal $341 = 11 \times 31$ tersusun: bilangan itu lolos uji
Fermat pada bilangan pokok $2$ padahal bukan \omterm{def:g12:arith:prime}{prima}. Pelajarannya: \omterm{def:g12:arith:congruence}{kekongruenan}
Fermat itu perlu, tetapi tidak cukup --- pengujian keprimaan
menuntut perkakas yang lebih tajam (dan memperolehnya, di jilid
universitas).

\textbf{13.} $3d \equiv 1 \pmod{20}$: jadi $d = 7$
(sebab $21 = 20 + 1$).

\textbf{14.} $c = 4^3 = 64 \equiv 31 \pmod{33}$.

\textbf{15.} $31 \equiv -2$: maka $(-2)^7 = -128$, dan
$-128 + 4 \times 33 = 4$: jadi teks sandinya terbaca kembali menjadi
$m = 4$. Kuncinya berputar.

\textbf{16.} Modulo $3$: jika $3 \nmid m$, maka
$m^2 \equiv 1$ (Fermat), sehingga
$m^{21} = m \cdot \left(m^2\right)^{10} \equiv m$; sedangkan jika
$3 \mid m$, kedua ruasnya $\equiv 0$. Modulo $11$:
$m^{10} \equiv 1$ atau $11 \mid m$, dan
$m^{21} = m \cdot \left(m^{10}\right)^2 \equiv m$. Jadi $3$
dan $11$ keduanya \omterm{def:g12:arith:divides}{membagi} $m^{21} - m$, dan karena \omterm{def:g12:arith:gcd}{saling prima} hasil kalinya
$33$ juga \omterm{def:g12:arith:divides}{membaginya} (Gauss): $m^{21} \equiv m \pmod{33}$. Eksponen
$ed = 21 = 1 + 20k$ dibangun justru agar kedua eksponen
Fermat ($2$ dan $10$, yang \omterm{def:g12:arith:divides}{membagi} $20$) lenyap.

\textbf{17.} Mengalikan dua \omterm{def:g12:arith:prime}{bilangan prima} sepanjang $300$ angka memakan waktu
sepersejuta detik; memperolehnya kembali dari hasil kalinya mengalahkan setiap
algoritme yang diketahui dan seluruh komputer di dunia --- kuncinya
jalan satu arah. (Bilangan $n = 33$ kita adalah jalan itu dalam ukuran mainan,
yang dapat ditempuh ke dua arah.)

\textbf{18.} Dengan menguji sisanya (atau membangun lewat Bézout):
$n \equiv 8 \pmod{15}$, jadi cacahnya $8, 23, 38, 53, \dots$
Ketunggalannya modulo $15$: dua penyelesaian berselisih kelipatan
$3$ sekaligus $5$, sehingga kelipatan $15$ (sebab $3$ dan $5$ \omterm{def:g12:arith:gcd}{saling prima}, Gauss).
Sang jenderal dengan $1000$ prajurit mengumumkan ``$8$'' lewat tiga
kali baris cepat --- muslihat pencacahan kepala dari zaman dahulu.

\textbf{19.} Dari $10 \equiv 1 \pmod 9$ diperoleh $10^k \equiv 1$, sehingga
$\sum d_k 10^k \equiv \sum d_k$: jadi sebuah bilangan dan jumlah angkanya
kongruen modulo $9$ (dan modulo $3$). Lalu
$10 \equiv -1 \pmod{11}$ memberi
$\sum d_k 10^k \equiv \sum (-1)^k d_k$: itulah kaidah berselang-selingnya.
Dua kaidah masa kanak-kanak, masing-masing satu baris.

\textbf{20.} \omterm{def:g12:arith:congruence}{Kekongruenan} mengubah sisa menjadi sebuah aritmetika
(Bagian I); Bézout mencetak balikan yang menyelesaikan \omterm{def:g12:arith:congruence}{kekongruenan}
linear dan $d$ milik RSA (pertanyaan 4, 13); teorema kecil
Fermat membuka dan menutup kuncinya (pertanyaan 15--16); perekatan
\omterm{def:g12:complex:modulus}{modulus} yang \omterm{def:g12:arith:gcd}{saling prima} mencacah prajurit dan menuntaskan buktinya
(pertanyaan 16, 18). Vonis atas Hardy: teorema paling murni yang
dikenalnya kini menjaga setiap pembelian --- kemurnian, jika diberi waktu,
adalah hal yang paling dapat diterapkan.

\end{solution}
